%!TEX root = ../main.tex

\section{Discussion}
\label{sec:discussion}

\subsection{Why calibrated far-field uncertainty should reduce heading error}
\noindent The guidance line in a detect-then-fit pipeline is anchored by the same
detections that are least reliable. Under a forward-facing field camera, rows near the
top of the image span few pixels, are partially occluded by foreground canopy, and project
many ground metres into a single image row; the detector's box edges there are
correspondingly diffuse. In our parameterisation $x = a\,y + b$, the heading
$\theta = \operatorname{atan}(a)$ is the controller-relevant readout, and the slope $a$ is
governed disproportionately by the spread of the centroids in $y$. A handful of far-field
centroids that are displaced by tens of pixels can therefore swing $\theta$ by several
degrees even when every near-field box is perfectly placed --- which is precisely the regime
an equal-weight least-squares fit cannot defend against, because it treats a noisy far box
and a sharp near box as equally informative.

CALM-Row addresses this at the only point where the information is still available: the
distributional DFL head~\cite{gfl} already emits, per box edge $e\in\{l,t,r,b\}$, a
softmax distribution $p_e(k)$ over $R=\mathit{reg\_max}$ bins, whose variance
$\sigma_e^2 = \sum_k (k-\mu_e)^2\,p_e(k)$ widens exactly when the edge is uncertain. By
linear error propagation through the centroid, the per-box horizontal variance is
$\sigma_{c_x}^2 = (s^2/4)(\sigma_l^2+\sigma_r^2)$, and the precision weight
$w_i = 1/\sigma_{c_x,i}^2$ down-weights diffuse far boxes inside the GLS normal equations.
Because the weighting acts multiplicatively on each residual, its leverage is largest
exactly where $\sigma_{c_x}^2$ is largest --- the far field --- and negligible where the
detector is confident. This is the mechanistic reason we expect the largest A/B gap to
appear on the held-out far rows rather than on the full-frame average, and it is why the
pre-registered far-field analysis (Section~\ref{sec:results}) is the discriminating test
rather than the headline mean. An illustrative synthetic check (Section~\ref{sec:res-synth}) only confirms that the mechanism behaves as derived; the empirical claim rests entirely on the real, un-manipulated far-field test.

Crucially, the weights are not free parameters to be tuned but a \emph{calibrated}
quantity. The distance-aware calibration loss (DDC) is a Gaussian negative-log-likelihood
that ties $\sigma_{c_x}^2$ to the realised squared centroid error on matched anchors, so a
predicted variance of a given magnitude must, on validation data, correspond to localisation
errors of the matching magnitude. This is what lets us report reliability (ECE, coverage of
realised heading error at $68/95\%$) rather than assert that the network has learned
``depth''; the uncertainty is a measured, checkable property (reported as ECE and coverage).

\subsection{Decomposing the gain: weighting alone (B0) versus training (B)}
\noindent A reviewer's natural question is whether any improvement comes from the
\emph{idea} of uncertainty weighting or from the \emph{extra training}. Our design answers
this by construction. Arm~B0 reads out a precision-weighted GLS line from the raw DFL
variance of a \emph{stock} detector trained with standard losses --- it shares one
checkpoint with the equal-weight baseline~A and the RANSAC baseline~A$'$, so the three are
compared paired per frame with zero training-induced variance between them. Arm~B is a
separate checkpoint trained with the DDC and LSL auxiliary losses. The contrast
B0 versus A isolates the value of the \emph{free} variance the head already produces; the
contrast B versus B0 isolates what the calibration-and-line-aware training \emph{adds} on
top. We report both regardless of which dominates: if B0 already beats A, the result is a
near-zero-cost readout change requiring no retraining; if the gain only materialises in B,
the training is doing the work and we say so plainly. Either outcome is an honest, useful
finding, and the decomposition forecloses the ``the gain is just more compute'' critique.

The expected qualitative pattern follows from the two losses. Raw DFL variance is correlated
with localisation error but is not guaranteed to be \emph{calibrated} in magnitude, so B0
should capture the gross near/far ordering but may mis-scale the weights; DDC re-scales the
variance to match realised error, and LSL --- whose gradients flow through $w_i$ into the
DFL logits, with the far look-ahead row explicitly included in its RMSE term --- additionally
teaches the network to be \emph{decisive about which far boxes to distrust}, pushing
unreliable far edges toward flatter distributions. We therefore expect
B0 $\le$ B on the far-field heading metric, but we treat the size of the
B$-$B0 increment as an empirical question to be filled from the results
(Section~\ref{sec:results}) rather than asserted.

\subsection{What CALM-Row is, and is not}
\noindent We are deliberate about the scope of the contribution. CALM-Row is \emph{not} an
architectural novelty: it introduces no new neck, no new backbone block, and --- in its
training losses --- no new learnable parameters, since those losses are computed entirely
from the variance the DFL head~\cite{gfl} already produces. The detector is a stock
YOLOv8n / YOLOv10n, and the exported INT8/RKNN or TensorRT graph is a vanilla YOLO; CALM-Row
runs as post-NMS CPU arithmetic and adds zero NPU operators. This is a feature, not a
hedge: it means any improvement is attributable to the estimator and the navigation-task
training, not to an unaccounted increase in capacity, and it cleanly separates our claim
from feature-fusion necks such as ASFF~\cite{asff}, which improve the features a detector computes rather than the line fit.

Positioned honestly, CALM-Row is a \emph{navigation-native, calibrated-uncertainty
estimator} that sits between a detector and a controller. Detection-uncertainty methods such
as KL-Loss~\cite{klloss} and Gaussian-YOLO~\cite{gaussianyolo} model box uncertainty to
improve \emph{detection}; we instead read a calibrated per-detection centroid covariance
out of the existing DFL distribution and propagate it into the downstream geometric estimate
that the application actually consumes. Relative to crop-row pipelines --- equal-weight
least-squares centerlines on YOLOX maize rows~\cite{diao_yolox}, segmentation-plus-scan
methods on the CRDLD benchmark~\cite{cth1,cth2}, polynomial regressors~\cite{rowdetr}, and
row-anchor regressors ported from lane detection~\cite{alnet} --- the differentiator is not
a new output paradigm but that the line is fit with \emph{principled, calibrated weights}
and is trained and evaluated in navigation units (heading in degrees, cross-track in
centimetres) rather than detection mAP. We make no claim that CALM-Row's detector is more
accurate than these systems in mAP; the claim is confined to the guidance-line readout and
its calibration, evaluated under a controlled A/B in which only the line-fit axis varies.

\subsection{Limitations}
\label{sec:limitations}
\noindent We state the boundaries of the evidence explicitly, since several of them bound
how far the conclusions generalise.

\paragraph{Generalisation and its limits.}
Generalisation is assessed in two ways: in-domain on the held-out CRDLD test split, and as
\emph{cross-dataset} transfer to CRBD~\cite{crbd} (a different camera, crop and resolution),
evaluated without retraining. The CRDLD download used here carries no per-field metadata, so a
within-dataset leave-one-field-out protocol is unavailable; the cross-dataset CRBD transfer
takes its place and is, if anything, a stronger probe. The unit of generalisation is the
dataset, not the seed.

\paragraph{Flat-ground homography for metric units.}
Centimetre-scale metrics require an image-to-ground homography, which we recover under a
\emph{flat-ground} assumption (or anchor to known inter-row spacing where intrinsics are
unavailable). On sloped or uneven terrain this assumption is violated, and a non-zero camera
pitch inflates the metric error most at the far look-ahead distance --- the same region
CALM-Row targets. We therefore report degrees (which are controller-independent and free of
the homography) as the primary metric, give centimetre figures only where the calibration is
defensible, and provide an explicit pitch-to-centimetre error budget so the metric claims
are not over-read.

\paragraph{Cross-track is a static geometric proxy, not a closed-loop result.}
We report cross-track as the lateral offset (px, and cm where the homography permits)
between the predicted and ground-truth guidance lines at the look-ahead row. We deliberately
do \emph{not} run a closed-loop pure-pursuit simulation and do \emph{not} report an
intervention rate: both require sequential field video and a controller model and are left to
future work. The controller-independent heading error is our primary navigation metric, and a
real-robot validation likewise remains future work.

\paragraph{Dependence on the DFL spread.}
CALM-Row's weights are only as informative as the variance the DFL head carries. If
$R=\mathit{reg\_max}$ is too small, or the head collapses toward near-deterministic
distributions, the per-edge variance saturates and the weights lose discriminative power
between near and far boxes; the method would then degenerate toward equal-weight
least-squares. Stock Ultralytics uses $R=16$ for all model sizes, which we expect to carry
sufficient spread, and we ablate $R\in\{8,16\}$ to test this dependence rather than assume
it. The approach also presupposes a distributional regression head; it does not transfer
unchanged to detectors with a single-point box regressor.

\paragraph{Calibration after INT8 quantisation.}
The variance the head emits can shift under INT8 quantisation, so calibration established in
floating point need not survive on the deployed graph. We therefore re-fit the optional
post-hoc \texttt{CalibScale} affine on a small validation set after quantisation and report
calibration before versus after INT8. We treat ``calibration survives quantisation'' as a
hypothesis to be demonstrated, not assumed; if the post-INT8 reliability degrades, the
deployment-time recommendation is to re-calibrate, and the weighting --- being monotone in
the variance --- still preserves the near/far ordering even when the absolute scale drifts.

\paragraph{Ground-truth line extraction choices.}
The evaluation depends on a ground-truth guidance line, and its definition is a modelling
choice. For evaluation we use a \emph{mask-derived} centerline (a frozen row-wise robust
regression, Section~\ref{sec:setup-gtline}) that is independent of the detector, deliberately
avoiding the circularity of fitting the ground-truth line through ground-truth box centroids
when the very claim under test is weighted-versus-unweighted fitting through box centroids.
We cross-check the automatically extracted line against a human-annotated subset and report
inter-annotator agreement in degrees and centimetres. Residual sensitivity to the
extraction algorithm --- gap and occlusion handling, the choice of the navigation row, and
the robust-regression estimator --- is a limitation of any line-level benchmark and bounds
the resolution at which differences smaller than the annotation agreement can be trusted. For
training, the line-stability loss instead uses a line fit through \emph{ground-truth box
centroids}, which is detector-independent and hence non-circular with respect to the
evaluation; the two definitions serve different purposes and are reported separately.